\scanpage{87}
\mspage{43}
disques troués autour du germe de
variété $D=(\overline s)_{\mathrm{red}}$ en $s\dots$, et l'extension
de $N_{s'}$ par $\pi$ contenue dans l'ext. $E_{U_K}$ de
$E_K$ par $\pi_1$.
\margin{\deleted{\unsure{c'est l'extension} $E_{U_{\unsure{\widetilde K|K}}}$ \unsure{correspondante|correspondant}} \unsure{\deleted{\unsure{à le $s$|à la $s$}}}}
La section donnée (plutôt,
\unsure{construit|construite} via $f_K$) de l'ext. $E_{U_K}$, en
définit une de la sous extension. \deleted{\unsure{Si on|Comme}}

L'interprétation géométrique de ces constructions
\deleted{\unsure{considérer|considère} d'autre part les scindages \unsure{de|des}}
profinies est \unsure{la|le} \unsure{suivante|suivant} : Soit
$\deleted{\unsure{(\widetilde S)_w|(\widetilde S)_s|\widetilde Z}}\inserted{\unsure{\widetilde{\underline O}_s|\underline O_s}}$ le hensélisé (pas strict) de $\underline O_{S,s}=\underline O_s$,
et $\widetilde S=\operatorname{Spec}(\widetilde{\underline O}_s)$ \inserted{(trait hensélien)}, $\widetilde K$ son corps des
fonctions, on a $U_{\widetilde S}$ sur $\widetilde S$, $U_{\widetilde K}$ sur $\widetilde K$,
et \deleted{l'extension}
\[
(23)\quad N_{s'}\quad
\begin{gathered}[t]
\text{(groupe de}\\
\text{décomposition}\\
\text{de }E_{K,\overline K}\text{ en }s')
\end{gathered}
\simeq E_{\widetilde K}
\]
et \unsure{la|le} \deleted{\unsure{suite|seule|sous}}\inserted{$E$} extension de ce groupe par $\pi$
n'est autre que
\diagram{\begin{tikzcd}[ampersand replacement=\&,column sep=large,row sep=1.2em]
(24) \& 1 \arrow[r,dash] \& \pi \arrow[r] \& E_{U_{\widetilde K}} \arrow[r] \arrow[ddd] \& E_{\widetilde K} \arrow[r] \arrow[ddd,bend right=18] \& 1 \\
\& \& \begin{gathered}\Vert\\\pi_1(U_{\overline K})\end{gathered} \& \& \begin{gathered}\Vert\\\operatorname{Gal}(\overline K/\widetilde K)\end{gathered} \& \\
\text{s'envoyant dans} \& \& \& \& \& \\
(25) \& 1 \arrow[r,dash] \& \pi \arrow[r] \& E_U \arrow[r] \& E_{\widetilde O} \arrow[r] \& 1 \\
\& \& \& \& \begin{gathered}\text{\tikz[baseline=(rel.base)]{\node[rotate=90,inner sep=0pt] (rel) {$\simeq$};}}\\\operatorname{Gal}(\overline k/k)\end{gathered} \&
\end{tikzcd}}
($k=k(s)$ est le corps résiduel).
\margin{86}
